\documentclass[11pt]{article}

\usepackage[margin=1in]{geometry}
\usepackage{amsmath,amssymb,amsthm}
\usepackage{mathtools}

\newtheorem{lemma}{Lemma}

\title{The Covariance Estimator Used in the Current Implementation}
\author{}
\date{}

\begin{document}

\maketitle

\section{Construction}

Fix a fold $m\in[M]$. All nuisance quantities below are estimated using only
$\mathcal D^{(-m)}$. The group-specific mean estimators are denoted by
$\mu_X^{(-m)}$ and $\mu_Z^{(-m)}$. In the current implementation, these are
either the sample means or coordinatewise hard-thresholded sample means.
The observations in the two groups are centered separately using these
estimators and then pooled.

Let $K$ be the number of estimated principal components. The vectors
$v_1^{(-m)},\ldots,v_K^{(-m)}$ are obtained from the pooled training sample.
For sparse PCA, the current implementation uses \texttt{PMA::SPC}; for the
hard-thresholded alternative, it first calculates the ordinary principal
components and then thresholds their coordinates. Each resulting vector is
normalized to have Euclidean norm one.

For $k\in[K]$, the corresponding eigenvalue is estimated by the empirical
variance in the direction $v_k^{(-m)}$:
\begin{align}
\lambda_k^{(-m)}
={}&\frac{1}{|\mathcal D^{(-m)}|}
\left[
\sum_{X_i\in\mathcal D^{(-m)}}
  \left\{(X_i-\mu_X^{(-m)})^\top v_k^{(-m)}\right\}^2
\right.\nonumber\\
&\left.\hspace{35mm}+
\sum_{Z_i\in\mathcal D^{(-m)}}
  \left\{(Z_i-\mu_Z^{(-m)})^\top v_k^{(-m)}\right\}^2
\right].
\label{eq:estimated-eigenvalue}
\end{align}
The residual noise variance is estimated by distributing the unexplained
empirical covariance trace equally over the remaining $p-K$ directions:
\begin{align}
\widehat\sigma^{2(-m)}
=\frac{1}{p-K}\Bigg[
&\frac{1}{|\mathcal D^{(-m)}|}
\left\{
\sum_{X_i\in\mathcal D^{(-m)}}
\|X_i-\mu_X^{(-m)}\|^2
+
\sum_{Z_i\in\mathcal D^{(-m)}}
\|Z_i-\mu_Z^{(-m)}\|^2
\right\}
-\sum_{k=1}^K\lambda_k^{(-m)}
\Bigg].
\label{eq:estimated-noise}
\end{align}

Equations~\eqref{eq:estimated-eigenvalue} and
\eqref{eq:estimated-noise} implicitly define the spiked covariance estimator
used by the current implementation:
\begin{equation}
\Sigma^{(-m)}
=
\sum_{k=1}^K
\left\{\lambda_k^{(-m)}-\widehat\sigma^{2(-m)}\right\}
v_k^{(-m)}v_k^{(-m)\top}
+\widehat\sigma^{2(-m)}I_p.
\label{eq:current-covariance-estimator}
\end{equation}
The code does not need to form the $p\times p$ matrix in
\eqref{eq:current-covariance-estimator}. Instead, it uses this spectral form
directly when calculating
\begin{equation*}
\left(\lambda_k^{(-m)}I_p-\Sigma^{(-m)}\right)^+.
\end{equation*}

\section{The leading-component case}

When $K=1$, \eqref{eq:current-covariance-estimator} becomes
\begin{equation}
\Sigma^{(-m)}
=
\left\{\lambda_1^{(-m)}-\widehat\sigma^{2(-m)}\right\}
v_1^{(-m)}v_1^{(-m)\top}
+\widehat\sigma^{2(-m)}I_p.
\label{eq:rank-one-estimator}
\end{equation}
Because $\|v_1^{(-m)}\|=1$, this construction gives the exact identity
\begin{equation}
\Sigma^{(-m)}v_1^{(-m)}
=\lambda_1^{(-m)}v_1^{(-m)}.
\label{eq:estimated-eigenpair}
\end{equation}
Moreover, \eqref{eq:rank-one-estimator} is positive semidefinite whenever
$\widehat\sigma^{2(-m)}\geq0$ and
$\lambda_1^{(-m)}\geq\widehat\sigma^{2(-m)}$.
Thus, in the leading-component case, the estimated eigenpair used in the
influence-function calculation is an exact eigenpair of the implicitly
constructed covariance estimator.

\section{Reduction of the operator-norm requirement}

The following lemma records sufficient conditions under which the estimator
in \eqref{eq:current-covariance-estimator} satisfies the operator-norm
condition used in the one-step theorem.

\begin{lemma}[Operator-norm consistency of the spiked estimator]
\label{lem:spiked-covariance-rate}
Suppose $K$ is fixed and
\begin{equation}
\Sigma
=
\sum_{k=1}^K(\lambda_k-\sigma^2)v_kv_k^\top
+\sigma^2I_p+R^{[1]},
\label{eq:population-spiked-decomposition}
\end{equation}
where $v_1,\ldots,v_K$ are orthonormal, the $\lambda_k$ and $\sigma^2$ are
uniformly bounded, and $\|R^{[1]}\|=o(n^{-1/4})$. Suppose also that
$v_1^{(-m)},\ldots,v_K^{(-m)}$ are orthonormal and, after sign alignment,
\begin{equation*}
\max_{k\in[K]}|\lambda_k^{(-m)}-\lambda_k|=o_P(n^{-1/4}),
\qquad
|\widehat\sigma^{2(-m)}-\sigma^2|=o_P(n^{-1/4}),
\end{equation*}
and
\begin{equation*}
\max_{k\in[K]}\|v_k^{(-m)}-v_k\|=o_P(n^{-1/4}).
\end{equation*}
Then the estimator in \eqref{eq:current-covariance-estimator} satisfies
\begin{equation*}
\|\Sigma^{(-m)}-\Sigma\|=o_P(n^{-1/4}).
\end{equation*}
\end{lemma}

\begin{proof}
Subtracting \eqref{eq:population-spiked-decomposition} from
\eqref{eq:current-covariance-estimator} and applying the triangle inequality
gives
\begin{align*}
\|\Sigma^{(-m)}-\Sigma\|
\leq{}&
\sum_{k=1}^K
\left|\lambda_k^{(-m)}-\widehat\sigma^{2(-m)}\right|
\left\|v_k^{(-m)}v_k^{(-m)\top}-v_kv_k^\top\right\|\\
&+
\sum_{k=1}^K
\left|
\lambda_k^{(-m)}-\lambda_k
-\widehat\sigma^{2(-m)}+\sigma^2
\right|
+|\widehat\sigma^{2(-m)}-\sigma^2|
+\|R^{[1]}\|.
\end{align*}
For unit vectors with aligned signs,
\begin{equation*}
\left\|v_k^{(-m)}v_k^{(-m)\top}-v_kv_k^\top\right\|
\leq 2\|v_k^{(-m)}-v_k\|.
\end{equation*}
The estimated eigenvalues and residual noise variance are bounded in
probability under the stated assumptions. Since $K$ is fixed, substituting
the assumed rates into the preceding bound proves the result.
\end{proof}

Lemma~\ref{lem:spiked-covariance-rate} separates the required covariance rate
into four tasks: controlling the spiked-model residual $R^{[1]}$, estimating
the leading eigenvalues, estimating the leading eigenvectors, and estimating
the residual noise variance. The implementation computes the corresponding
estimators, but their rates require assumptions on the covariance structure,
the sparsity of the leading eigenvectors, the eigengaps, and the growth of
$p$ relative to the training sample size.

\section{Multiple estimated components}

For $K>1$, the interpretation of
\eqref{eq:current-covariance-estimator} as a spectral decomposition requires
the estimated vectors to be orthonormal. Under this condition,
\begin{equation*}
\Sigma^{(-m)}v_k^{(-m)}
=\lambda_k^{(-m)}v_k^{(-m)},
\qquad k\in[K].
\end{equation*}
The current call to \texttt{PMA::SPC} uses its default
\texttt{orth = FALSE}. Consequently, when several sparse components are
estimated, their loading vectors need not be exactly orthogonal. In that
case, \eqref{eq:current-covariance-estimator} remains a well-defined symmetric
matrix, but the displayed eigenpair identity need not hold. This issue does
not arise for $K=1$. For $K>1$, the theoretical spectral calculation therefore
requires either an estimator that returns orthonormal loading vectors or a
separate analysis of the non-orthogonal factor representation.

\end{document}
